% BEGIN VOL07-EXPANSION VII02
\section{Supplementary worked examples}
\label{sec:vii02-pedagogy}

\begin{example}[Cartesian and polar coordinates]\label{ex:vii02-ped-01}
On the punctured plane, Cartesian coordinates \((x,y)\) and polar coordinates \((r,\theta)\) overlap after choosing an angular branch.  The transition map
\[
(r,\theta)\longmapsto (r\cos\theta,r\sin\theta)
\]
and its local inverse are smooth.  Smooth compatibility is therefore a statement about transition maps, not about the names of coordinates.
\end{example}

\begin{example}[Stereographic smooth structure on the sphere]\label{ex:vii02-ped-02}
The north- and south-pole stereographic charts on \(S^n\) have overlap map \(u\mapsto u/\lVert u\rVert^2\) on \(\mathbb R^n\setminus\{0\}\).  This map is smooth with smooth inverse, so the two-chart topological atlas generates the standard smooth structure on \(S^n\).
\end{example}

\begin{example}[Restricting a smooth atlas]\label{ex:vii02-ped-03}
If \(M\) is smooth and \(U\subset M\) is open, restricting every chart to \(U\) produces a smooth atlas on \(U\).  The transition functions are restrictions of the original smooth transition functions.  Thus open submanifolds inherit smooth structure without any new compatibility proof.
\end{example}

\section{Graded supplementary exercises}
The following exercises add a deliberate progression from concrete calculation to structural proof, hypothesis testing, application, and synthesis.

\subsection*{Standard computations and constructions}

\begin{exercise}[Polar transition]\label{exr:vii02-ped-01}
Compute the Jacobian of \((r,\theta)\mapsto(r\cos\theta,r\sin\theta)\).
\end{exercise}
\begin{hint}
Differentiate with respect to \(r\) and \(\theta\).
\end{hint}
\begin{solution}
The Jacobian matrix is \(\begin{pmatrix}\cos\theta&-r\sin\theta\\ \sin\theta&r\cos\theta\end{pmatrix}\), with determinant \(r\), nonzero on the punctured-plane domain \(r>0\).
\end{solution}

\begin{exercise}[Sphere overlap]\label{exr:vii02-ped-02}
Verify that \(u\mapsto u/\lVert u\rVert^2\) is its own inverse on \(\mathbb R^n\setminus\{0\}\).
\end{exercise}
\begin{hint}
Apply the map twice and simplify the norm.
\end{hint}
\begin{solution}
If \(v=u/\|u\|^2\), then \(\|v\|^2=1/\|u\|^2\), hence \(v/\|v\|^2=u\).
\end{solution}

\begin{exercise}[Restricted atlas]\label{exr:vii02-ped-03}
Describe an atlas for an open annulus \(1<\|x\|<2\subset\mathbb R^2\).
\end{exercise}
\begin{hint}
The identity map already gives a global smooth chart.
\end{hint}
\begin{solution}
The annulus is itself open in \(\mathbb R^2\), so the identity chart is a smooth global chart.
\end{solution}

\begin{exercise}[Cylinder coordinates]\label{exr:vii02-ped-04}
Give smooth local coordinates on \(S^1\times\mathbb R\).
\end{exercise}
\begin{hint}
Use any smooth angular chart on \(S^1\) and the identity on \(\mathbb R\).
\end{hint}
\begin{solution}
The product of the two coordinates gives a smooth chart with values in an open subset of \(\mathbb R^2\).
\end{solution}

\begin{exercise}[Coordinate change]\label{exr:vii02-ped-05}
On \(x>0\), compute the transition from Cartesian coordinates to \((u,v)=(\log x,y/x)\).
\end{exercise}
\begin{hint}
Solve for \(x\) and \(y\).
\end{hint}
\begin{solution}
The inverse is \(x=e^u,\ y=e^u v\), which is smooth; hence the coordinates are smoothly compatible on \(x>0\).
\end{solution}

\subsection*{Proofs}

\begin{exercise}[Compatibility is transitive]\label{exr:vii02-ped-06}
Prove that smooth compatibility of charts is transitive.
\end{exercise}
\begin{hint}
Write the transition from the first chart to the third as a composition through the second.
\end{hint}
\begin{solution}
On common overlap components, \(\chi\circ\varphi^{-1}=(\chi\circ\psi^{-1})\circ(\psi\circ\varphi^{-1})\); a composition of smooth maps is smooth, and similarly for the inverse.
\end{solution}

\begin{exercise}[Maximal atlas]\label{exr:vii02-ped-07}
Show that every smooth atlas is contained in a unique maximal smooth atlas.
\end{exercise}
\begin{hint}
Adjoin all charts smoothly compatible with the given atlas.
\end{hint}
\begin{solution}
The compatible charts are mutually compatible by transitivity and form an atlas.  Any maximal atlas containing the original one must contain all of them, proving uniqueness.
\end{solution}

\begin{exercise}[Open restriction]\label{exr:vii02-ped-08}
Prove that an open subset of a smooth \(n\)-manifold inherits a canonical smooth structure.
\end{exercise}
\begin{hint}
Restrict every chart domain to the open subset.
\end{hint}
\begin{solution}
The restrictions cover the open subset and their transition maps are restrictions of the original smooth transitions.
\end{solution}

\begin{exercise}[Product smooth structure]\label{exr:vii02-ped-09}
Prove that products of smooth manifolds carry a natural smooth product atlas.
\end{exercise}
\begin{hint}
Take products of compatible charts.
\end{hint}
\begin{solution}
Transition functions split as products of smooth transition functions, hence are smooth maps between open subsets of Euclidean space.
\end{solution}

\subsection*{Counterexamples and hypothesis tests}

\begin{exercise}[Topological compatibility is enough]\label{exr:vii02-ped-10}
Test the claim that two topological charts automatically belong to the same smooth atlas.
\end{exercise}
\begin{hint}
Compare coordinates \(x\) and \(x^{1/3}\) on \(\mathbb R\).
\end{hint}
\begin{solution}
False.  The transition \(x\mapsto x^{1/3}\) is not differentiable at \(0\), so topological compatibility does not imply smooth compatibility.
\end{solution}

\begin{exercise}[One smooth chart determines every smooth structure]\label{exr:vii02-ped-11}
Test whether a single local chart fixes the smooth structure on a whole disconnected manifold.
\end{exercise}
\begin{hint}
Change charts on another component.
\end{hint}
\begin{solution}
False.  A chart on one component carries no compatibility information about charts on a disjoint component.
\end{solution}

\begin{exercise}[Smoothness is coordinate-free without checks]\label{exr:vii02-ped-12}
Test the claim that a formula looking smooth in one arbitrary topological chart defines a smooth object.
\end{exercise}
\begin{hint}
The chart itself may not belong to the smooth atlas.
\end{hint}
\begin{solution}
False.  Smoothness must be tested in charts from the chosen smooth structure; an arbitrary homeomorphism can destroy differentiability.
\end{solution}

\subsection*{Applications and investigations}

\begin{exercise}[Smooth torus]\label{exr:vii02-ped-13}
Explain how the product atlas gives \(T^2=S^1\times S^1\) its standard smooth structure.
\end{exercise}
\begin{hint}
Take products of smooth circle charts.
\end{hint}
\begin{solution}
Each product chart lands in \(\mathbb R^2\), and overlap maps are products of one-dimensional smooth transitions.
\end{solution}

\begin{exercise}[Smooth quotient by translations]\label{exr:vii02-ped-14}
Explain why \(\mathbb R^n/\mathbb Z^n\) inherits a smooth atlas.
\end{exercise}
\begin{hint}
Use small Euclidean balls on which the quotient is injective.
\end{hint}
\begin{solution}
Local inverses of the quotient supply coordinates; transition maps are restrictions of integer translations, hence smooth.
\end{solution}

\subsection*{Challenge problems}

\begin{exercise}[Equivalent atlases]\label{exr:vii02-ped-15}
Show that two atlases define the same smooth structure exactly when their union is a smooth atlas.
\end{exercise}
\begin{hint}
Translate equality of maximal atlases into pairwise compatibility.
\end{hint}
\begin{solution}
If the maximal structures agree, every chart in one is compatible with every chart in the other, so the union is smooth.  Conversely, a smooth union is contained in one maximal atlas, forcing the generated maximal atlases to coincide.
\end{solution}

\begin{exercise}[Pulling back a smooth structure]\label{exr:vii02-ped-16}
Let \(h:X\to M\) be a homeomorphism and \(M\) a smooth manifold.  Construct a smooth structure on \(X\) making \(h\) a diffeomorphism.
\end{exercise}
\begin{hint}
Pull every chart \((U,\varphi)\) back to \((h^{-1}U,\varphi\circ h)\).
\end{hint}
\begin{solution}
The pulled-back transition functions are exactly those on \(M\), hence smooth; by construction the chart representations of \(h\) are identities.
\end{solution}

% END VOL07-EXPANSION VII02
